\documentclass[11pt,a4paper]{article}
\usepackage[utf8]{inputenc}
\usepackage[T1]{fontenc}
\usepackage[brazil]{babel}
\usepackage{lmodern}
\usepackage{amsmath,amssymb,amsthm}
\usepackage{booktabs}
\usepackage{float}
\usepackage[margin=2.5cm]{geometry}
\usepackage{pgfplots}
\pgfplotsset{compat=1.16}
\usepackage[hidelinks]{hyperref}

\newtheorem{teorema}{Teorema}
\newtheorem{proposicao}[teorema]{Proposição}
\newtheorem{lema}[teorema]{Lema}
\theoremstyle{remark}
\newtheorem{observacao}[teorema]{Observação}

\title{O perímetro da elipse via formas modulares:\\ funções teta, o nome de Jacobi e a série de Eisenstein \texorpdfstring{$E_2$}{E2}}
\author{}
\date{Novembro de 2026}

\begin{document}
\maketitle

\begin{abstract}
Usando a teoria das formas modulares, escrevemos o perímetro da elipse de forma
exata em termos da função teta $\theta_3$, da série de Eisenstein de peso 2,
$E_2$, e do nome de Jacobi $q$. Há duas versões, ligadas pela transformação
modular $\tau\mapsto-1/\tau$. Na versão para elipses achatadas, a transformação
modular faz surgir naturalmente o logaritmo que descreve a singularidade do
perímetro no caso degenerado. Como o nome é sempre $q\le e^{-\pi}\approx0{,}043$
no ramo usado, poucas parcelas das $q$-séries bastam. O resultado são fórmulas elementares,
só com raízes e logaritmo. A mais simples tem erro relativo máximo de $4{,}8\times10^{-6}$
em todo o intervalo; uma versão ligeiramente maior chega a $5\times10^{-12}$. Ambas
superam muito a segunda aproximação de Ramanujan (pior caso $4{,}0\times10^{-4}$).
\end{abstract}

\section{As fórmulas}

Sejam $a\ge b\ge0$ os semieixos e $r=b/a$. Use o ramo A se $r\ge1/\sqrt2$ e o
ramo B se $r<1/\sqrt2$.

\paragraph{Versão simples.}
\begin{align}
\textbf{Ramo A:}\quad &\lambda=\frac{\sqrt a-\sqrt b}{2(\sqrt a+\sqrt b)},\quad
T=(1+2\lambda)^2,\quad
\boxed{P\approx\frac{2\pi a}{3}\left[T\,(1+r^2)+\frac{1-24\lambda^2-72\lambda^4}{T}\right]}
\label{eq:A}\\[6pt]
\textbf{Ramo B:}\quad &e=\sqrt{1-r^2},\quad \mu=\frac{1-\sqrt e}{2(1+\sqrt e)},\quad T=(1+2\mu)^2,\notag\\
&\boxed{P\approx4a\left[\frac1T+\frac{\ln(1/\mu)}{6}\left(T\,(1+r^2)-\frac{1-24\mu^2-72\mu^4}{T}\right)\right]}
\label{eq:B}
\end{align}
Erro relativo máximo: $4{,}8\times10^{-6}$, atingido perto de $r=1/\sqrt2$.

\paragraph{Versão precisa.} É a mesma fórmula, trocando, em cada ramo, $\lambda$ (ou $\mu$) por
\begin{equation}\label{eq:precisa}
q=\lambda+2\lambda^5,\qquad
T=(1+2q+2q^4)^2,\qquad
E_2\text{-termo}=1-24\,(q^2+3q^4+4q^6+7q^8),
\end{equation}
e usando $\ln(1/q)$ no ramo B. Erro relativo máximo: $5\times10^{-12}$.

Para calcular $\mu$ quando $b$ é muito pequeno, sem cancelamento numérico, use
$\mu=\dfrac{r^2}{2(1+e)(1+\sqrt e)^2}$.

\section{Dedução}

\subsection{Objetos modulares}

Seja $k'=b/a$ o módulo complementar e $k=e=\sqrt{1-k'^2}$ a excentricidade. Sejam
$K=K(k)$ e $E=E(k)$ as integrais elípticas completas, e $K'=K(k')$, $E'=E(k')$. O
perímetro é $P=4aE$. Definimos:
\begin{itemize}
\item o \textbf{nome de Jacobi} $q=e^{-\pi K'/K}\in[0,1)$;
\item as \textbf{funções teta}
$\theta_2(q)=2\sum_{n\ge0}q^{(n+1/2)^2}$,
$\theta_3(q)=1+2\sum_{n\ge1}q^{n^2}$ e
$\theta_4(q)=1+2\sum_{n\ge1}(-1)^nq^{n^2}$;
\item a \textbf{série de Eisenstein de peso 2}
$E_2(q)=1-24\sum_{n\ge1}\sigma_1(n)\,q^n$, em que $\sigma_1(n)$ é a soma dos
divisores de $n$. Os primeiros valores são $\sigma_1=1,3,4,7,6,12,\dots$.
\end{itemize}

\begin{lema}[Jacobi]\label{lema:jacobi}
\[
K=\frac\pi2\,\theta_3(q)^2,\qquad \sqrt{k'}=\frac{\theta_4(q)}{\theta_3(q)},\qquad
\frac{K}{K'}=\frac{\pi}{\ln(1/q)} .
\]
\end{lema}

\subsection{O nome a partir de \texorpdfstring{$a$}{a} e \texorpdfstring{$b$}{b}}

\begin{proposicao}\label{prop:nome}
Seja $\lambda=\dfrac{1-\sqrt{k'}}{2(1+\sqrt{k'})}=\dfrac{\sqrt a-\sqrt b}{2(\sqrt a+\sqrt b)}$. Então
\[
\lambda=\frac{\theta_2(q^4)}{2\,\theta_3(q^4)}
=q\,\frac{1+q^8+q^{24}+\cdots}{1+2q^4+2q^{16}+\cdots},
\]
e, invertendo a série,
\[
q=\lambda+2\lambda^5+15\lambda^9+150\lambda^{13}+1707\lambda^{17}+\cdots
\]
\end{proposicao}
\begin{proof}
Pelo Lema~\ref{lema:jacobi},
$\lambda=\frac{\theta_3-\theta_4}{2(\theta_3+\theta_4)}$. Usando as identidades
$\theta_3(q)-\theta_4(q)=2\theta_2(q^4)$ e $\theta_3(q)+\theta_4(q)=2\theta_3(q^4)$, que saem
separando os termos pares e ímpares das séries, obtém-se a igualdade. A inversão é a série
clássica de Jacobi e Weierstrass para o nome.
\end{proof}

Note que $\lambda$ é \emph{elementar}: depende só de $\sqrt a$ e $\sqrt b$.
Para $k'\ge1/\sqrt2$ vale $q\le e^{-\pi}\approx0{,}0432$, e a série converge muito
rapidamente.

\subsection{\texorpdfstring{$E$}{E} em termos de \texorpdfstring{$E_2$}{E2}}

\begin{proposicao}[Ramanujan]\label{prop:E2}
\begin{equation}\label{eq:EK}
\frac EK=\frac{1+k'^2}{3}+\frac{E_2(q^2)}{3\,\theta_3(q)^4}.
\end{equation}
\end{proposicao}
\begin{proof}[Comentário]
Esta é uma reformulação de uma fórmula de Ramanujan para a série $P(q^2)=E_2(q^2)$:
$P(q^2)=z^2\big(3\tfrac EK-2+k^2\big)$, com $z=\tfrac{2K}{\pi}=\theta_3^2$ (ver
B.~C.~Berndt, \emph{Ramanujan's Notebooks}, Parte III, cap.~17). Como
$2-k^2=1+k'^2$, ela equivale a \eqref{eq:EK}. Conferimos a identidade numericamente com
38 dígitos. No círculo ($q=0$), ela dá $\frac23+\frac13=1$.
\end{proof}

\subsection{Ramo A: elipses moderadas}

Multiplicando \eqref{eq:EK} por $K=\frac\pi2\theta_3^2$ e por $4a$, obtemos a identidade
\emph{exata}
\begin{equation}\label{eq:Aexata}
P=\frac{2\pi a}{3}\left[\theta_3(q)^2\,(1+k'^2)+\frac{E_2(q^2)}{\theta_3(q)^2}\right].
\end{equation}
Truncar $q\approx\lambda$, $\theta_3\approx1+2q$ e $E_2(q^2)\approx1-24q^2-72q^4$ dá \eqref{eq:A}.

\subsection{Ramo B: transformação modular e o logaritmo}

Quando $b\to0$, $q\to1$ e as séries de \eqref{eq:Aexata} deixam de convergir bem. A
transformação modular $\tau\mapsto-1/\tau$ troca os papéis de $k$ e $k'$. Seja
$p=e^{-\pi K/K'}$ o nome \emph{complementar}, que tende a $0$ quando $b\to0$. Pelos
Lemas e Proposições acima aplicados ao módulo $k'$:
\[
K'=\frac\pi2\theta_3(p)^2,\qquad
\frac{E'}{K'}=\frac{1+k^2}{3}+\frac{E_2(p^2)}{3\theta_3(p)^4},\qquad
\frac{K}{K'}=\frac{\ln(1/p)}{\pi},
\]
e $p$ se obtém de $\mu=\frac{1-\sqrt k}{2(1+\sqrt k)}$ pela mesma série da
Proposição~\ref{prop:nome}.

A \textbf{relação de Legendre}, $EK'+E'K-KK'=\frac\pi2$, dá
\[
E=\frac{\pi}{2K'}+K\Big(1-\frac{E'}{K'}\Big)
=\frac{\pi}{2K'}+K\left[\frac{1+k'^2}{3}-\frac{E_2(p^2)}{3\theta_3(p)^4}\right].
\]
Substituindo $K'$ e $K=K'\ln(1/p)/\pi$:
\begin{equation}\label{eq:Bexata}
P=4a\left\{\frac{1}{\theta_3(p)^2}+\frac{\ln(1/p)}{6}\left[\theta_3(p)^2(1+k'^2)-\frac{E_2(p^2)}{\theta_3(p)^2}\right]\right\}.
\end{equation}
Esta identidade também é \emph{exata}. Truncando, obtém-se \eqref{eq:B}.

\begin{observacao}[a origem do logaritmo]
Nas fórmulas racionais ou algébricas, a dificuldade está na singularidade logarítmica do
perímetro quando $b\to0$. Aqui ela aparece sozinha: é o fator
$\ln(1/p)=\pi K/K'$ produzido pela transformação modular. Quando $b\to0$,
temos $p\approx\mu\approx r^2/16$ e a chave em \eqref{eq:Bexata} se comporta como $\tfrac{r^2}{2}\ln\tfrac{4}{r}$, o que
reproduz a expansão clássica $P\approx4a\big[1+\frac{r^2}{2}(\ln\frac4r-\frac12)+\cdots\big]$.
\end{observacao}

\subsection{Escolha do ramo e erro de truncamento}

Os dois nomes satisfazem $\ln q\cdot\ln p=\pi^2$, e são iguais quando $k=k'=1/\sqrt2$, com
$q=p=e^{-\pi}$. Escolhendo sempre o ramo de nome menor, temos
$q\le e^{-\pi}=0{,}0432$. Os primeiros termos desprezados são da ordem de:
\begin{itemize}
\item versão simples: $2\lambda^5$ no nome e $2q^4$ em $\theta_3$, que dão $\sim10^{-6}$;
\item versão precisa: $15\lambda^9$, $2q^9$ e $24\cdot6\,q^{10}$, que dão $\sim10^{-12}$.
\end{itemize}
Isso concorda com os erros medidos (Seção~\ref{sec:num}). O pior caso é sempre a fronteira
$r=1/\sqrt2$, e o erro cai rapidamente longe dela.

\begin{teorema}
As expressões \eqref{eq:Aexata} e \eqref{eq:Bexata}, com as séries completas, são exatas. As
versões truncadas \eqref{eq:A}–\eqref{eq:B} e \eqref{eq:precisa}:
\begin{enumerate}
\item são exatas no círculo ($\lambda=0$) e no segmento ($\mu=0$, em que o termo com $\ln$ tende a $0$);
\item usam apenas operações elementares (raízes, logaritmo e potências);
\item têm erro relativo máximo de $4{,}8\times10^{-6}$ (simples) e $5\times10^{-12}$ (precisa).
\end{enumerate}
\end{teorema}

\section{Verificação numérica}\label{sec:num}

Comparação com $P=4aE(e)$ calculado com 40 dígitos (\texttt{mpmath}), em 4000 valores
de $b/a$ e em valores extremos até $b/a=10^{-37}$.

\begin{table}[H]
\centering
\begin{tabular}{cccc}
\toprule
$b/a$ & simples & precisa & Ramanujan II\\
\midrule
0.99 & $3.3\times10^{-12}$ & $1.6\times10^{-25}$ & $2.4\times10^{-28}$\\
0.9  & $4.0\times10^{-8}$  & $2.1\times10^{-16}$ & $3.8\times10^{-18}$\\
0.8  & $8.1\times10^{-7}$  & $1.3\times10^{-13}$ & $6.7\times10^{-15}$\\
0.71 & $4.6\times10^{-6}$  & $4.6\times10^{-12}$ & $4.7\times10^{-13}$\\
0.7  & $4.3\times10^{-6}$  & $2.9\times10^{-12}$ & $7.0\times10^{-13}$\\
0.5  & $1.9\times10^{-7}$  & $2.0\times10^{-15}$ & $4.6\times10^{-10}$\\
0.3  & $3.5\times10^{-9}$  & $2.8\times10^{-20}$ & $7.4\times10^{-8}$\\
0.2  & $1.7\times10^{-10}$ & $3.8\times10^{-23}$ & $8.5\times10^{-7}$\\
0.1  & $9.0\times10^{-13}$ & $5.2\times10^{-28}$ & $1.2\times10^{-5}$\\
0.01 & $1.8\times10^{-20}$ & $<10^{-38}$ & $2.4\times10^{-4}$\\
0.001& $2.8\times10^{-28}$ & $<10^{-38}$ & $3.8\times10^{-4}$\\
\midrule
máximo & $4.8\times10^{-6}$ & $5.0\times10^{-12}$ & $4.0\times10^{-4}$\\
\bottomrule
\end{tabular}
\caption{Erro relativo (em módulo). Perto do círculo, Ramanujan II é mais preciso: supera a
versão precisa para $b/a\gtrsim0.68$ e a simples para $b/a\gtrsim0.37$. Em todo o resto, as fórmulas modulares vencem com folga, e o pior caso
delas é muito menor.}
\end{table}

\begin{figure}[H]
\centering
\begin{tikzpicture}
\begin{axis}[
  width=0.85\textwidth, height=7cm,
  xlabel={$b/a$}, ylabel={erro relativo $|\,\cdot\,|$},
  ymode=log, ymin=1e-36, ymax=1e-2, xmin=0, xmax=1,
  legend pos=south east, grid=major]
\addplot[blue, thick] table[x=x, y=simples] {erros.dat};
\addlegendentry{simples}
\addplot[green!50!black, thick] table[x=x, y=precisa] {erros.dat};
\addlegendentry{precisa}
\addplot[red, thick, dashed] table[x=x, y=ram2] {erros.dat};
\addlegendentry{Ramanujan II}
\addplot[black, densely dotted] coordinates {(0.7071,1e-36) (0.7071,1e-2)};
\end{axis}
\end{tikzpicture}
\caption{Erro relativo em função de $b/a$. A linha pontilhada marca a troca de ramo
($b/a=1/\sqrt2$), onde o erro é máximo.}
\end{figure}

\section{Observações}

\begin{itemize}
\item \textbf{Área matemática.} A construção usa ferramentas da teoria dos números e das
formas modulares: funções teta de Jacobi, a série de Eisenstein $E_2$ (com a função soma de
divisores $\sigma_1$), o nome $q$ e a transformação modular $\tau\mapsto-1/\tau$.
\item \textbf{Uso prático.} Só há raízes, potências e um logaritmo, então a fórmula funciona em qualquer
calculadora científica.
\item \textbf{Ineditismo.} As identidades usadas são clássicas (Jacobi, Legendre e Ramanujan),
e o cálculo numérico de integrais elípticas por funções teta é conhecido. Não encontramos,
porém, a fórmula do perímetro na forma de dois ramos, elementar e truncada, apresentada
aqui. A busca (novembro de 2026) não foi exaustiva. O mais honesto é apresentá-la como uma
aplicação nova de resultados clássicos, e não como um teorema novo.
\end{itemize}

\end{document}
